\documentclass[11pt]{article}
\usepackage[a4paper,margin=26mm]{geometry}
\usepackage{amsmath,amssymb,amsthm,mathtools}
\usepackage[T1]{fontenc}
\usepackage{lmodern}
\usepackage[colorlinks=true,allcolors=blue]{hyperref}
\newtheorem{theorem}{Theorem}
\newtheorem{lemma}{Lemma}
\newcommand{\Prob}{\mathop{\rm Pr}}
\newcommand{\F}{\mathbb F}
\newcommand{\N}{\mathbb N}
\newcommand{\Z}{\mathbb Z}
\newcommand{\Aut}{\operatorname{Aut}}
\newcommand{\id}{\operatorname{id}}
\begin{document}
\begin{center}
{\large\bfseries A degree distinguisher for SQIsign hint distributions}\\[3pt]
{\small Exact finite theorem and accompanying Lean definitions}
\end{center}
\vspace{-0.6em}

The attack distinguishes the stationary-target and degree-weighted hint laws in
\emph{A Complete Security Proof of SQIsign}~\cite{proof}, in the supplied revision dated
August 1, 2025. It reads one public isogeny degree and tests whether it is small.
The statement below fixes the raw-morphism counting convention and explicitly
completes sampling steps that otherwise may fail.

\paragraph{Curves, maps, and scalar stripping.}
Let $p\ge49$ be prime with $p\equiv3\pmod4$, put $k=\F_{p^2}$, and write
\[
 E_A:\ y^2=x^3+Ax^2+x,\qquad
 \mathcal M_p=\{A\in k:A^2\ne4,\ E_A\text{ supersingular}\}.
\]
An isogeny is an actual pointed algebraic morphism over $k$, with its target
coefficient $A$ retained. Distinct morphisms to the same equation are counted
separately; neither post-isomorphism classes nor representation strings are sampled.
Set
\begin{align*}
e&=\min\{r\in\N:p\le4^r\}, &
B&=2^e=2^{\lceil\log_2\sqrt p\rceil},\\
W(d)&=\prod_{\ell^a\parallel d}(\ell+1)^a, &
Z(B)&=\sum_{d=1}^B W(d).
\end{align*}
In particular $p<B^2<4p$ and $B<p$. All maps of degree at most $B$ are separable.

If $\varphi:X\to Y$ has degree $d\le B$ and kernel $K$, its scalar content is the
largest positive divisor $m$ of $d$ for which $X[m]\subseteq K$. There is a unique
isogeny $\chi:X\to Y$ with
\[
 \varphi=\chi\circ[m]_X=[m]_Y\circ\chi,\qquad
 \ker\chi=[m]_X K,\qquad C:=\deg\chi=d/m^2,
\]
and $\ker\chi$ is cyclic. Thus stripping preserves the actual target, not just
its geometric $j$-invariant.

\paragraph{The two total finite experiments.}
Fix supersingular normalized Montgomery curves $X_U,X_S$ over $k$; they need not
be equal. Define the optional degree observations $U$ and $S$ as follows.
\begin{itemize}
\item $U$: draw $d\in\{1,\ldots,B\}$ with probability $W(d)/Z(B)$, then draw uniformly
from all raw degree-$d$ isogenies $X_U\to E_A$, $A\in\mathcal M_p$.
Compute the cyclic-component degree $C$. Return $\bot$ only if $C=B$, the
zero-complementary-degree case. Otherwise return $C$; in particular, retain
noncyclic raw degree-$B$ maps whose cyclic component has degree $C<B$.
\item $S$: draw $A$ uniformly from $\mathcal M_p$, then uniformly select a raw
isogeny $X_S\to E_A$ of degree $<B$. An empty connecting fiber returns $\bot$
without resampling $A$. Otherwise return its cyclic-component degree $C$.
\end{itemize}
Both experiments keep their full original normalizers. In particular failure
is not followed by conditioning on success. Replacing only the $C=B$ failure
by any other output cannot decrease acceptance; the bound below also covers
such non-resampling completions retaining every valid $C<B$ output. For a successful degree
$C=2^n c$ with $c$ odd, the complementary degree in Experiment~3 and Algorithm~5
is $B/2^n-c$: it is a positive integer because $C<B$. The subsequent cyclic
isogeny can therefore be sampled from a nonempty finite set.

\begin{theorem}[Finite raw-morphism distinguishing theorem]\label{thm:finite}
Under these conventions, let $s\ge784$ be an integer satisfying $8s^2\le B$, and set
\[
 L=\left\lfloor\frac{B}{4s^2}\right\rfloor,\qquad
 A_{B,s}(C)=\begin{cases}1&C\in\N\text{ and }C\le L,\\0&C=\bot\text{ or }C>L.\end{cases}
\]
Then, for every pair $X_U,X_S$ as above,
\begin{equation}\label{eq:main}
 \Prob[A_{B,s}(U)=1]-\Prob[A_{B,s}(S)=1]
 \ \ge\ \frac1{1024s^3}.
\end{equation}
The same signed bound holds after averaging over any common finite public-context
law that chooses $X_U$ and $X_S$. It also holds after arbitrary randomized
continuations and valid encodings that preserve the observed optional degree.
For any transcript of $q\ge1$ hints, the test uses only the first hint. Later
challenges and hints may be sampled with arbitrary conditional randomness
preserving that first observation; failures remain per-hint failure symbols.

Suppose the exposed component has an efficient representation in the sense of
Definition~2.11 of~\cite{proof}: a polynomial-time public validator returns its
actual degree, and representation lengths are polynomially bounded in the public
parameter lengths. Then the test, including degree extraction, runs in polynomial
time. Algorithm~6 preserves the degree of the first exposed component, so the same
bound applies after that push whenever it is performed on the successful hint.
\end{theorem}

The source curves in Experiment~4 have different roles: $X_U$ is the public-key
curve before pushing, while $X_S$ is the challenge codomain in the simulated
experiment. The theorem's two-source formulation is what permits averaging over
the public key and challenge. The event-probability advantage is the absolute
difference of acceptance probabilities;~\eqref{eq:main} proves its sign as well.
Success above $1/2$ in the corresponding fair-bit guessing game is half this gap.

\paragraph{Geometry and the exact sampling multiplicities.}
For $d<p$, the actual elliptic torsion is
$X[d]\cong(\Z/d\Z)^2$. Finite-subgroup classification by Hermite normal form (HNF)
therefore gives exactly
\[
 \sigma_1(d)=\sum_{a\mid d}a
\]
order-$d$ kernels. Write an HNF as $(a,e,b)$ with $a,e>0$, $0\le b<a$, and
$ae=d$. Its content is $m=\gcd(a,e,b)$, and division by $m$ gives a primitive
HNF, hence a cyclic subgroup. In the elliptic point group the stripped subgroup
is $[m]K$, of order $d/m^2$. Moreover $X[r]\subseteq K$ if and only if $r\mid m$
for positive $r\mid d$; for any positive $r$ prime to $p$, full torsion
contained in $K$ necessarily has $r^2\mid d$. Actual isogeny factorization through $[m]$, and surjectivity of
$[m]$ on geometric points, identify this subgroup with $\ker\chi$.

A supersingular geometric class has respectively $2,3,6$ normalized Montgomery
parameters when $j=0$, $j=1728$, or otherwise. Its automorphism group has
respectively $6,4,2$ elements. Each fixed kernel consequently has exactly $12$
raw quotient morphisms when all Montgomery targets are counted. For a fixed
target equation, the number is $|\Aut(E_A)|$, independently of the kernel.
Thus raw uniform sampling induces uniform kernel sampling, both degree by degree
and within each nonempty target fiber.

These are also the counts over $k$, not merely over $\overline{\F}_p$.
Geometric supersingular normalized Montgomery parameters lie in $k$; for
$p\equiv3\pmod4$ their $p^2$-Frobenius is $[-p]$, and they have $(p+1)^2$ rational
points. Any geometric isogeny between these models commutes with their common
scalar Frobenius and descends uniquely to $k$. One descent proof expresses a
function as $u(x)+yv(x)$: coefficient-Frobenius invariance fixes the normalized
numerators and denominators of $u,v$, whose coefficients then lie in $k$.
Consequently commutation with Frobenius gives Galois-equivariant pullbacks and
actual morphism descent.

The Hasse polynomial $H(T)=\sum_i\binom{(p-1)/2}{i}^{\!2}T^i$ has simple roots
and describes the supersingular Legendre locus. Its fiber counts, or equivalently
the Montgomery parameter counts above, give the stationary geometric-class law
\begin{equation}\label{eq:stationary}
 \mu(j)=\frac{24}{(p-1)|\Aut(E_j)|}\le\frac{12}{p-1}.
\end{equation}
This is precisely the $j$-marginal of uniform sampling from $\mathcal M_p$.
The descent and scalar-factorization statements are essential: equality of
$j$-invariants alone would not identify the raw finite-field samplers.

\begin{lemma}[Weight and normalizer bounds]\label{lem:weight}
For every $d\ge1$ and $B\ge1$,
$W(d)\ge\sigma_1(d)$ and $Z(B)\le4B^2$.
\end{lemma}
\begin{proof}
At a prime power the binomial theorem gives
$(\ell+1)^a\ge1+\ell+\cdots+\ell^a$; multiply over primes.
Let $g$ be multiplicative, with $g(1)=1$ and $g(\ell^a)=(\ell+1)^{a-1}$ for
$a\ge1$. The identity
\[
 (\ell+1)^a=\ell^a+\sum_{j=1}^a\ell^{a-j}(\ell+1)^{j-1}
\]
gives $W=\id*g$. Hence
\[
 Z(B)=\sum_{n\le B}g(n)\sum_{k\le B/n}k
 \le B^2\sum_{n\le B}\frac{g(n)}{n^2}.
\]
The nonnegative Euler factors are
$\sum_{a\ge0}g(\ell^a)/\ell^{2a}=1+1/(\ell^2-\ell-1)$.
The factor at $2$ is $2$, and for $n\ge3$,
\[
 1+\frac1{n^2-n-1}\le\frac{(n-1)^2}{n(n-2)}.
\]
Bounding prime factors by all integer factors and telescoping gives
$2\prod_{n=3}^B (n-1)^2/[n(n-2)]=4(B-1)/B<4$ for $B\ge3$.
The cases $B=1,2$ are immediate.
\end{proof}

\begin{proof}[Proof of Theorem~\ref{thm:finite}]
More generally, let $L,M\ge1$ and $M^2L<B$, and use $U\le L$ or $S\le L$
to denote a nonfailure degree at most $L$. Acceptance in the target-first
experiment forces its target to be the codomain of a cyclic map of degree at most
$L$. Scalar stripping preserves that target, so at most
\[
 \sum_{c=1}^L\sigma_1(c)=\sum_{a=1}^L a\left\lfloor\frac La\right\rfloor\le L^2
\]
target classes can be accepted. Equation~\eqref{eq:stationary} yields
$\Prob[S\le L]\le12L^2/(p-1)$.
This upper bound even permits an arbitrary conditional connecting-map law.

For $U$, consider the distinct HNFs
\[
 (a,e,b)=(mc,m,mt),\qquad
 1\le m\le M,\quad1\le c\le L,\quad0\le t<c.
\]
Each has degree $m^2c<B$, content exactly $m$, and cyclic-component degree $c$;
they remain distinct when their original degrees coincide. The $c$ kernels in a
fixed $(m,c)$ category contribute probability
\[
 \frac{W(m^2c)}{Z(B)}\frac{c}{\sigma_1(m^2c)}
 \ge\frac{c}{4B^2}.
\]
The raw morphism multiplicity $12$ cancels. Summation gives
\[
 \Prob[U\le L]\ge\frac{ML(L+1)}{8B^2},\qquad
 \Prob[U\le L]-\Prob[S\le L]
 \ge\left(\frac M8-49\right)\frac{L^2}{B^2},
\]
since $12B^2/(p-1)<48p/(p-1)\le49$. Every favorable category has raw degree
$d<B$. Discarding all raw degree-$B$ outputs is therefore a conservative lower
bound; restoring its valid noncyclic outputs can only increase acceptance.
The full normalizer proves the bound for failure only at $C=B$, and for any
further non-resampling completion of that sole invalid case.

Now take $M=s$ and $L=\lfloor B/(4s^2)\rfloor$. The hypotheses imply
$B/(8s^2)\le L\le B/(4s^2)$, $s^2L<B$, and $s/8-49\ge s/16$.
The last display is at least
$(s/16)(1/(64s^4))=1/(1024s^3)$.
The bound is uniform in both source curves, so averaging its signed form over a
common public context preserves it. Appending data or encoding a map does not
change any test on its preserved degree. Algorithm~6's final splitting retains
both exposed component degrees. Its odd-kernel push preserves cardinality and
cyclicity by the actual dual-isogeny identities and coprimality with the
power-of-two challenge.

For public extraction, run the validator of Definition~2.11 on the first exposed
component representation, reject failure, and test $4s^2C\le B$ on its binary
degree. This is exactly $C\le L$ and requires only polynomial-time arithmetic
in addition to the stipulated polynomial-time validator.
\end{proof}

\paragraph{Asymptotic consequence and scope.}
For a family with a fixed polynomial-time validator, polynomially bounded
representation lengths and $\log p$ in a security parameter $\lambda$, and
$B\ge8(\lambda+784)^2$ eventually, choose $s=\lambda+784$.
One hint gives inverse-polynomial advantage $1/[1024(\lambda+784)^3]$.
For example $p\ge2^\lambda$ eventually supplies the threshold condition.
Thus computational indistinguishability fails for these precisely specified
hint laws and valid public representations. This is a distinguishing attack,
not a forgery or an endomorphism-ring algorithm. It does not contradict a
reduction that explicitly retains a hint-distinguishing advantage term.
The statement does not replace either sampling law by uniform sampling of
representation strings, and it does not silently condition empty fibers or
endpoint failures away.

\paragraph{Lean correspondence.}
\begingroup\raggedright
The certificate's definitions retain actual point subgroups, V\'elu function-field
isogenies, normalized finite-field coefficients, and raw morphisms.
\texttt{ScalarStripGeometry} and \texttt{CyclicComponent} connect HNF content to
actual same-target factorization. \texttt{MontgomeryDescent},
\texttt{MontgomeryLegendre}, \texttt{IsogenyDescent}, and
\texttt{MontgomeryMorphismDescent} supply parameter descent, negative scalar
Frobenius, and morphism descent. \texttt{RawDegreeSampler},
\texttt{RawTargetSampler}, and \texttt{RawTargetGeometry} specify the raw laws,
empty fibers, and actual component observation; \texttt{RawSamplerAttack} gives
the signed interior lower bound and averaging over two source curves.
\texttt{SourceEndpoint} transfers it to the source law retaining all valid
noncyclic endpoint outcomes and failing exactly when $C=B$.
\texttt{RationalAttack} states the bound on literal finite-field map laws.
\texttt{PushedEncoding} and \texttt{TranscriptEncoding} connect actual pushes
and arbitrary fresh later hints to the first public observation.
\texttt{EndToEnd} and \texttt{EndToEndFamilies} provide the finite public
certificate and encoded-family nonnegligibility theorem.
\texttt{ResponseCutoff} identifies the integer cutoff with the printed ceiling,
and \texttt{PublicObservation} composes a polynomial degree decoder with the
verified binary comparison. Its work bound concerns an explicit charged
Boolean/list computation model; the degree decoder is a representation
interface, not an asserted implementation of an arbitrary external encoding.
The build and axiom-audit reports accompany the Lean sources; verification is
performed on the named declarations and their dependencies.
\par\endgroup

\begin{thebibliography}{1}
\bibitem{proof}
M.~A.~Aardal, A.~Basso, L.~De~Feo, S.~Patranabis, and B.~Wesolowski.
\emph{A Complete Security Proof of SQIsign}.
IACR ePrint 2025/379, supplied revision dated August 1, 2025.
\url{https://eprint.iacr.org/2025/379}.
\end{thebibliography}
\end{document}
